% ============================================================
%  Paper VIII  --  Working draft
%  Scale-Separated Dyadic Cancellation for PSWF Galerkin Operators
%  Status: working draft (revised v6: submission-ready final)
% ============================================================
\documentclass[12pt]{amsart}

% --- packages ---
\usepackage{amsmath,amssymb,amsthm}
\usepackage{mathtools}
\usepackage{hyperref}
\usepackage{enumerate}

% --- theorem environments ---
\newtheorem{theorem}{Theorem}[section]
\newtheorem{proposition}[theorem]{Proposition}
\newtheorem{conjecture}[theorem]{Conjecture}
\newtheorem{corollary}[theorem]{Corollary}
\newtheorem{lemma}[theorem]{Lemma}
\theoremstyle{definition}
\newtheorem{assumption}[theorem]{Assumption}
\newtheorem{definition}[theorem]{Definition}
\theoremstyle{remark}
\newtheorem{remark}[theorem]{Remark}

% --- macros ---
\newcommand{\norm}[1]{\lVert#1\rVert}
\newcommand{\abs}[1]{\lvert#1\rvert}
\newcommand{\lam}[1]{\lambda_{#1}^{(c)}}
\newcommand{\FAi}{F_{\mathrm{Ai}}}
\newcommand{\NSh}{N_{\mathrm{Sh}}}
\newcommand{\Bm}{\mathcal{B}_m}

% ============================================================
\begin{document}
% ============================================================

\title[Scale-Separated Dyadic Cancellation]
      {Scale-Separated Dyadic Cancellation\\
       for PSWF Galerkin Operators:\\
       Two-Scale Structure of the Spectral Transition}

\author{[Author]}
\date{\today}
\maketitle

% ------------------------------------------------------------
\begin{abstract}
We study the eigenvalue transition of the PSWF concentration
operator in the high-bandwidth regime and establish two
classes of results.

\emph{Rigorous.}
Under Assumption~\ref{ass:gap} (uniform spectral gap
lower bound) and Assumption~B-strong \cite{paperVI},
we verify hypotheses~(H1)--(H3) of Paper~VII \cite{paperVII}
via the \emph{Dyadic Separation Principle}:
for $j \in \mathcal{B}_m(i)$, the ratio
$|\lambda_j - \lambda_{j+1}|/|\lambda_i - \lambda_j|$
is bounded by $C \cdot 2^{-m}$ independently of~$c$,
because the factor $(c/2)^{-1/3}$ cancels between
numerator and denominator.
This cancellation is a scale-combinatorial mechanism
requiring only monotonicity and a uniform gap lower
bound, and is therefore independent of any detailed
spectral asymptotics beyond that uniform local gap
condition.

\emph{Empirical.}
We present numerical evidence for a universal scaling law
governing the global transition profile
(Conjecture~\ref{conj:ULW}):
\[
  \lambda_l \approx \tfrac{1}{2}\operatorname{erfc}
  \!\left(\frac{\ln 2\cdot(l-\NSh-\beta(c))}{S(c)}\right),
  \quad S(c) = (\log c/\pi)^{2/3},
\]
with cross-bandwidth residuals of order $2\cdot 10^{-2}$
across $c \in \{50,100,200\}$, decreasing with~$c$.
The centering parameter $\beta(c)$ drifts slowly
($\beta(50)\approx -0.42$, $\beta(200)\approx -0.36$);
its limiting behavior is an open problem.
\end{abstract}

% ============================================================
\section{Introduction and Context}
\label{sec:intro}
% ============================================================

Paper~VII \cite{paperVII} established an abstract dyadic
cancellation theorem (Theorem~2.1) and reduced the
contraction problem to three hypotheses (H1)--(H3).
H2 was stated as Conjecture~6.2 and is reduced here
to Assumption~\ref{ass:gap}.
This paper is a \emph{conditional reduction theorem}:
under Assumption~\ref{ass:gap} and B-strong,
hypotheses (H1)--(H3) are verified and the contraction
estimate follows.

\medskip\noindent
\textbf{Structural novelty.}
The novelty of the present paper is the identification
of the Dyadic Separation Principle as a scale-combinatorial
mechanism driven solely by monotonicity and a uniform
gap condition, independent of any detailed spectral
asymptotics beyond that local gap bound.
This is not a renaming of hypothesis~H2: it is the
observation that H2 requires no Airy approximation,
no erfc structure, and no specific spectral asymptotics
whatsoever — only monotonicity and a lower bound on
consecutive eigenvalue gaps at scale $(c/2)^{-1/3}$.

\medskip\noindent
\textbf{Three-scale structure.}
Our analysis reveals three separated layers:
\begin{itemize}
\item \emph{Combinatorial scale} (Dyadic Separation):
  requires only a uniform gap lower bound at scale
  $(c/2)^{-1/3}$; this is Assumption~\ref{ass:gap},
  Lemma~\ref{lem:F}, and Corollary~\ref{cor:C}.
  This is the rigorous core, independent of the
  specific form of the eigenvalue profile.
\item \emph{Local scale} (Airy / WKB):
  provides a refined description of increments via
  Airy functions; this is Proposition~\ref{prop:U}
  and Lemma~\ref{lem:A}.
  It supplies one route to verifying
  Assumption~\ref{ass:gap} and the upper bounds
  needed for Term~2 of Corollary~\ref{cor:C},
  but is not a logical prerequisite for the
  Dyadic Separation Principle.
\item \emph{Global scale} (Landau--Widom):
  the eigenvalue profile follows an erfc law with
  logarithmic width $S(c)$; this is
  Conjecture~\ref{conj:ULW}, supported by numerics.
\end{itemize}

\medskip\noindent
\textbf{The Dyadic Separation Principle.}
For $j \in \mathcal{B}_m(i)$:
both $|\lambda_j-\lambda_{j+1}|$ and $|\lambda_i-\lambda_j|$
scale as $c^{-1/3}$, but $|\lambda_i-\lambda_j|$ carries
an extra $2^m$ from $|i-j| \ge 2^m$, so the ratio is
$O(2^{-m})$ independently of $c$.

\medskip\noindent
\textbf{What is not proved here.}
Assumption~\ref{ass:gap}, Proposition~\ref{prop:U},
and B-strong remain open;
see Remarks~\ref{rem:gap_route}, \ref{rem:propU_route},
and~\ref{rem:open}.
Conjecture~\ref{conj:ULW} is supported by numerical
evidence but has no analytic proof at present.

% ============================================================
\section{Setup and Notation}
\label{sec:setup}
% ============================================================

Throughout, $c \ge 1$; $\lambda_k \in (0,1)$ ordered
decreasing. $\NSh := 2c/\pi$.

\medskip\noindent
\textbf{Airy zeros and $A_0$.}
All zeros of $\mathrm{Ai}$ are negative;
$a_1 \approx -2.3381$ is the largest \cite[Sec.~9.9]{DLMF}.
Fix $A_0 \in (0,|a_1|)$; then $[-A_0,A_0]$ is zero-free.

\medskip\noindent
\textbf{Compactness constants (depend only on $A_0$):}
\begin{align*}
  c_* &:= \min_{[-A_0,0]}\FAi > 0,\\
  c_{\min} &:= \min_{[-A_0,A_0]}\mathrm{Ai}^2 > 0,\quad
  C_{\max} := \max_{[-A_0,A_0]}\mathrm{Ai}^2 < \infty.
\end{align*}

\medskip\noindent
\textbf{Lipschitz constant:}
$L_{A_0} := \sup_{[-A_0,A_0]}|\frac{d}{dx}\mathrm{Ai}^2|
= 2\sup|\mathrm{Ai}\,\mathrm{Ai}'| < \infty.$

\medskip\noindent
\textbf{Monotonicity.}
$\FAi$ is strictly decreasing
($\FAi' = -\mathrm{Ai}^2 \le -c_{\min} < 0$ on $[-A_0,A_0]$).

\medskip\noindent
\textbf{All constants} depend only on $A_0$; none
depend on $c$, $i$, $j$, $m$.

\medskip\noindent
\textbf{Rescaled index:}
$x_l := (l-\NSh)(c/2)^{-1/3}$;
spacing $x_{l+1}-x_l = (c/2)^{-1/3}$ exactly.

\medskip\noindent
\textbf{LW scaling variable:}
$S(c) := (\log c/\pi)^{2/3}$,\quad
$Z_l := \ln 2 \cdot (l - \NSh - \beta(c))/S(c)$.

\medskip\noindent
\textbf{Notation:}
$A_{ij} := P_{ij}c^{1/2}$,\;
$P_{ij} := c|\lambda_i-\lambda_j|/
((1-\lambda_i)(1-\lambda_j))$;\;
$\Bm(i) := \{j: 2^m \le |i-j| < 2^{m+1}\}$.

\medskip\noindent
\textbf{Bound on $1-\lambda_j$.}
Under Proposition~\ref{prop:U}\textup{(U)}
with $C_1 c_0^{-2/3} \le c_*/2$:
\begin{equation}\label{eq:onemlam}
  \tfrac{1}{2}\FAi(x_j) \le 1-\lambda_j \le 2\FAi(x_j)
\end{equation}
uniformly in all transition-window indices.

% ============================================================
\section{Spectral Assumptions and Conjectures}
\label{sec:structure}
% ============================================================

\begin{assumption}[Uniform Gap Condition]
\label{ass:gap}
There exists $\kappa_0 > 0$ such that for all
sufficiently large $c$ and all $l$ with
$|l - \NSh| \le A_0 c^{1/3}$:
\begin{equation}\label{eq:gapbound}
  \lambda_l - \lambda_{l+1} \ge \kappa_0\,(c/2)^{-1/3}.
\end{equation}
No further structural information on the eigenvalue
profile is required for Lemma~\ref{lem:F}.
\end{assumption}

\begin{remark}[Routes to Assumption~\ref{ass:gap}]
\label{rem:gap_route}
Three routes to~\eqref{eq:gapbound} are known or plausible:
\begin{enumerate}[(i)]
\item Via Proposition~\ref{prop:U}(U): for sufficiently
  large $c$, the Airy approximation implies
  $\kappa_0 = \tfrac{1}{2}c_{\min}$.
\item Via first-order WKB: a monotone comparison
  argument at leading order, without Airy matching.
  This route does not require (U) and may be more
  accessible.
\item Heuristic support from numerical computation
  combined with asymptotic stability arguments
  (not a proof; included for orientation only).
\end{enumerate}
None of these routes is currently complete.
\end{remark}

\begin{proposition}[ORX Airy Approximation]
\label{prop:U}
There exist $C_1, C_2 > 0$ (depending only on $A_0$)
such that for all sufficiently large $c$ and all $l$
with $|l-\NSh| \le A_0 c^{1/3}$:
\begin{itemize}
\item[\textup{(U)}] $\abs{\lambda_l - \FAi(x_l)}
  \le C_1 c^{-2/3}$.
\item[\textup{(U')}] $\abs{(\lambda_l-\lambda_{l+1})
  -(\FAi(x_l)-\FAi(x_{l+1}))}
  \le C_2 c^{-2/3}$.
\end{itemize}
Both bounds are uniform in $l$.
Proposition~\ref{prop:U}\textup{(U)} implies
Assumption~\ref{ass:gap} with
$\kappa_0 = \tfrac{1}{2}c_{\min}$.
\end{proposition}

\begin{remark}[Scope of Proposition~\ref{prop:U}]
\label{rem:propU_scope}
Proposition~\ref{prop:U} is a \emph{local} statement
at the Airy scale $(c/2)^{-1/3}$.
It makes no claim about the global eigenvalue profile.
The global erfc structure is addressed in
Conjecture~\ref{conj:ULW}.
\end{remark}

\begin{remark}[Why (U') is necessary for Lemma~\ref{lem:A}]
\label{rem:propU_uprime}
Condition~(U) bounds $|E_l|$ and $|E_{l+1}|$
individually where $E_l := \lambda_l - \FAi(x_l)$,
but does not control $E_l - E_{l+1}$.
Condition~(U') directly supplies this discrete
derivative bound and is a genuinely stronger condition.
\end{remark}

\begin{remark}[Route to proof of Proposition~\ref{prop:U}]
\label{rem:propU_route}
ORX \cite{OsipovRokhlinXiao} provides~(U) pointwise.
A proof of (U') additionally requires:
\begin{enumerate}[(i)]
\item Uniform WKB control with $O(c^{-2/3})$ remainder.
\item Airy matching with uniform error; yields~(U).
\item Stability under discrete differentiation;
  yields~(U').
\end{enumerate}
Step~(iii) is the additional content beyond~(U).
\end{remark}

\begin{conjecture}[Landau--Widom Global Scaling Law]
\label{conj:ULW}
Let $S(c) = (\log c/\pi)^{2/3}$.
There exists $\beta(c) \in \mathbb{R}$ such that
for all $l$ in the transition region,
\[
  \lambda_l = \frac{1}{2}
  \operatorname{erfc}\!\left(
  \frac{\ln 2 \cdot (l - \NSh - \beta(c))}{S(c)}
  \right)
  + O(S(c)^{-1}),
\]
and the first differences satisfy
\[
  \lambda_l - \lambda_{l+1}
  \approx \frac{\ln 2}{\sqrt{\pi}\,S(c)}\,e^{-Z_l^2}.
\]
\end{conjecture}

\begin{remark}[Numerical evidence for Conjecture~\ref{conj:ULW}]
\label{rem:ULW_numerics}
For $c \in \{50,100,200\}$, Gauss--Legendre
discretization ($n=400$) yields cross-bandwidth residuals
\[
  \sup_{c,c'}\sup_{Z \in [-2.5,2.5]}
  |\lambda_c(Z) - \lambda_{c'}(Z)| \approx 2.7 \cdot 10^{-2},
\]
decreasing with $c$, and
\[
  \beta(50)\approx -0.424,\quad
  \beta(100)\approx -0.380,\quad
  \beta(200)\approx -0.359.
\]
Fitting suggests $\beta(c) \to \beta(\infty)
\in [-0.30,-0.21]$; no analytic expression is known.
\end{remark}

\begin{remark}[Scale separation]
\label{rem:scales}
Assumption~\ref{ass:gap} and Lemma~\ref{lem:F}
act at the combinatorial scale $c^{-1/3}$.
Proposition~\ref{prop:U} resolves the local Airy
scale $(c/2)^{-1/3}$.
Conjecture~\ref{conj:ULW} describes the global
Landau--Widom scale $S(c) \sim (\log c)^{2/3}$.
The ratio $S(c)/(c/2)^{-1/3}
\sim c^{1/3}(\log c)^{2/3} \to \infty$
confirms genuine scale separation.
\end{remark}

% ============================================================
\section{Lemma A: Uniform Airy Increment}
\label{sec:lemmaA}
% ============================================================

\begin{lemma}[Uniform Airy Increment]
\label{lem:A}
Under Proposition~\ref{prop:U}\textup{(U')},
for all sufficiently large $c$ and all $l$ with
$|l-\NSh| \le A_0 c^{1/3}$:
\[
  \lambda_l - \lambda_{l+1}
  = (c/2)^{-1/3}\,\mathrm{Ai}(x_l)^2 + R_l,
  \quad |R_l| \le C_R\,c^{-2/3},
\]
where $C_R := L_{A_0}\cdot 2^{2/3} + C_2$.
Consequently:
\begin{equation}\label{eq:incr_bound}
  |\lambda_l-\lambda_{l+1}|
  \le C_{\max}(c/2)^{-1/3} + C_R\,c^{-2/3},
\end{equation}
and in particular
\begin{equation}\label{eq:incr_simple}
  |\lambda_l-\lambda_{l+1}|
  \le C_I\,(c/2)^{-1/3},
  \quad C_I := C_{\max} + C_R\cdot 2^{-1/3}.
\end{equation}
\end{lemma}

\begin{proof}
Write $\lambda_l-\lambda_{l+1}
= [\FAi(x_l)-\FAi(x_{l+1})]
+ [(\lambda_l-\lambda_{l+1})-(\FAi(x_l)-\FAi(x_{l+1}))]$.
By MVT and $\FAi'=-\mathrm{Ai}^2$:
$\FAi(x_l)-\FAi(x_{l+1})
= \mathrm{Ai}(\xi_l)^2(c/2)^{-1/3}$,
with $|\mathrm{Ai}(\xi_l)^2-\mathrm{Ai}(x_l)^2|
\le L_{A_0}(c/2)^{-1/3}$.
By~(U'): remainder $\le C_2 c^{-2/3}$.
Combining gives $|R_l| \le C_R c^{-2/3}$;
\eqref{eq:incr_bound}--\eqref{eq:incr_simple} follow
since $c^{-2/3}\le 2^{-1/3}(c/2)^{-1/3}$.
\end{proof}

% ============================================================
\section{Lemma B: Airy Ratio Bound}
\label{sec:lemmaB}
% ============================================================

\begin{lemma}[Airy Ratio Bound]
\label{lem:B}
$\mathrm{Ai}(x)^2/\FAi(x) \le C_{\max}/c_*$
for all $x \in [-A_0,A_0]$.
\end{lemma}

\begin{proof}
Numerator $\le C_{\max}$, denominator $\ge c_*>0$.
\end{proof}

% ============================================================
\section{Lemma F: Dyadic Separation Principle}
\label{sec:lemmaF}
% ============================================================

\begin{lemma}[Dyadic Separation Principle]
\label{lem:F}
Under Assumption~\ref{ass:gap},
for all sufficiently large $c$ and all $i \ne j$
in the transition window:
\begin{equation}\label{eq:specsep}
  |\lambda_i-\lambda_j|
  \ge \kappa_0\,(c/2)^{-1/3}\,|i-j|.
\end{equation}
\end{lemma}

\begin{proof}
WLOG $i < j$. By strict monotonicity and telescoping:
\[
  |\lambda_i - \lambda_j|
  = \sum_{l=i}^{j-1}(\lambda_l - \lambda_{l+1})
  \ge (j-i)\,\kappa_0\,(c/2)^{-1/3}.\qquad\qedhere
\]
\end{proof}

\begin{remark}[Full independence of Lemma~\ref{lem:F}]
\label{rem:Frobust}
The proof uses only Assumption~\ref{ass:gap}:
strict monotonicity and a gap lower bound at scale
$(c/2)^{-1/3}$.
No Airy functions, no erfc structure, no specific
form of $\lambda_l - \lambda_{l+1}$ is required.
Any mechanism verifying~\eqref{eq:gapbound} suffices.
\end{remark}

% ============================================================
\section{Corollary C: H2 in Truncated Dyadic Sense}
\label{sec:H2}
% ============================================================

\begin{corollary}[H2 in truncated dyadic sense]
\label{cor:C}
Fix $m_0\ge 1$.
Assume:
\begin{enumerate}[\textup{(a)}]
\item Assumption~\ref{ass:gap} (gap lower bound);
\item An upper bound
  $|\lambda_l - \lambda_{l+1}| \le C_I\,(c/2)^{-1/3}$
  for all $l$ in the transition window.
\end{enumerate}
Define
\[
  \delta_m
  := \frac{2C_I}{\kappa_0}\,2^{-m}
  + \frac{2C_{\max}}{c_*}(c/2)^{-1/3}
  + \frac{2C_R}{c_*}c^{-2/3}.
\]
For all sufficiently large $c$, all $i \le N$,
$m \ge m_0$, and $j,j+1 \in \Bm(i)$:
\[
  \frac{|A_{ij}-A_{i,j+1}|}{A_{ij}}
  \le \delta_m = O(2^{-m})+O(c^{-1/3}).
\]
For any $\varepsilon>0$: choose $m_0$ large, then $c$
sufficiently large; then $\sup_{m\ge m_0}\delta_m<\varepsilon$.

\smallskip\noindent
\emph{Condition~\textup{(b)} is implied by
Proposition~\ref{prop:U}\textup{(U')} via
\eqref{eq:incr_simple} of Lemma~\ref{lem:A}.
The bound on $1-\lambda_j$ used in Term~2 follows
from Proposition~\ref{prop:U}\textup{(U)}
via~\eqref{eq:onemlam}.}
\end{corollary}

\begin{proof}
\[
\frac{|A_{ij}-A_{i,j+1}|}{A_{ij}}
\lesssim
\underbrace{\frac{|\lambda_j-\lambda_{j+1}|}
            {|\lambda_i-\lambda_j|}}_{\text{Term 1}}
+
\underbrace{\frac{|\lambda_j-\lambda_{j+1}|}
            {1-\lambda_j}}_{\text{Term 2}}.
\]
\textbf{Term 1.}
Numerator $\le C_I(c/2)^{-1/3}$ by~(b).
Denominator $\ge \kappa_0(c/2)^{-1/3}\cdot 2^m$
by Lemma~\ref{lem:F}.
The factor $(c/2)^{-1/3}$ \emph{cancels exactly}:
\[
  \text{Term 1} \le \frac{C_I}{\kappa_0}\,2^{-m},
  \quad\text{independent of }c.
\]
\textbf{Term 2.}
By~\eqref{eq:incr_bound} and~\eqref{eq:onemlam}:
\[
  \frac{|\lambda_j-\lambda_{j+1}|}{1-\lambda_j}
  \le \frac{2C_{\max}}{c_*}(c/2)^{-1/3}
  + \frac{2C_R}{c_*}c^{-2/3} = O(c^{-1/3}).
\]
\end{proof}

\begin{remark}[The operative cancellation]
\label{rem:dyadic}
Both $|\lambda_j-\lambda_{j+1}|$ and $|\lambda_i-\lambda_j|$
are $O(c^{-1/3})$, but the denominator carries
the combinatorial factor $2^m$ from $|i-j|\ge 2^m$.
The $(c/2)^{-1/3}$ cancels, leaving $O(2^{-m})$
independent of~$c$.
The $O(c^{-1/3})$ term in $\delta_m$ comes from
Term~2 only and vanishes as $c\to\infty$.
\end{remark}

% ============================================================
\section{Centering Parameter $\beta(c)$}
\label{sec:beta}
% ============================================================

\begin{remark}[Empirical behavior of $\beta(c)$]
\label{rem:beta}
The parameter $\beta(c)$ in Conjecture~\ref{conj:ULW}
is determined numerically by median alignment.
Gauss--Legendre discretization ($n=400$) gives:
\[
  \beta(50)\approx -0.424,\quad
  \beta(100)\approx -0.380,\quad
  \beta(200)\approx -0.359.
\]
Fitting suggests $\beta(c)\sim A+B/\log(c)$ with
limit $\beta(\infty)\in[-0.30,-0.21]$,
but the model sensitivity is high.
The observed drift is consistent with a
discretization-induced centering effect,
but no structural identification of $\beta(c)$
in terms of spectral or arithmetic data is known.
\end{remark}

% ============================================================
\section{Lemma E: Core--WKB Split}
\label{sec:split}
% ============================================================

\begin{lemma}[Core--WKB Split]
\label{lem:E}
Fix $m_0 \ge 1$; assume Assumption~\ref{ass:gap},
Proposition~\ref{prop:U} (both \textup{(U)} and
\textup{(U')}), and B-strong \cite{paperVI}.
Then $\sum_{j\ne i}|T_{ij}| \le C(m_0,A,C_2)$
uniformly in $i$, $N$, and all sufficiently large $c$.
\end{lemma}

\begin{proof}
\textbf{Core ($m<m_0$).}
$|T_{ij}| \le C_2/|i-j|$ by B-strong
\cite[Lemma~5.2]{paperVII};
sum $\le 2C_2 m_0$.

\textbf{WKB ($m\ge m_0$).}
Corollary~\ref{cor:C} gives H2 via~(a) and~(b);
apply Theorem~2.1 of \cite{paperVII} with H1 and H3.
\end{proof}

% ============================================================
\section{Conditional Reduction}
\label{sec:closure}
% ============================================================

\begin{corollary}[Conditional Reduction]
\label{cor:closure}
Under Assumption~\ref{ass:gap},
Proposition~\ref{prop:U} (both \textup{(U)} and
\textup{(U')}), and B-strong,
hypotheses (H1)--(H3) of Theorem~2.1 in \cite{paperVII}
are verified, giving
$\sup_{i\le N}\sum_{j\ne i}|P_{ij}|
\le Cc^{-1/2}\log c$
and $\norm{D\mathcal{T}_c^{(N)}} < 1$
for all sufficiently large $c$.
\end{corollary}

\begin{remark}[Dependency chain]
\label{rem:deps}
\begin{align*}
  &\text{Ass.~\ref{ass:gap}}
    \Rightarrow\text{Lem.~\ref{lem:F}}
    \Rightarrow\text{Cor.~\ref{cor:C}\,(a)};\\
  &\text{Prop.~\ref{prop:U}(U')}
    \Rightarrow\text{Lem.~\ref{lem:A}}
    \Rightarrow\text{Cor.~\ref{cor:C}\,(b)};\\
  &\text{Prop.~\ref{prop:U}(U)}
    \Rightarrow\text{bound on }(1-\lambda_j)
    \Rightarrow\text{Cor.~\ref{cor:C}\,Term 2};\\
  &\text{Prop.~\ref{prop:U}(U)}
    \Rightarrow\text{Ass.~\ref{ass:gap} (one route)};\\
  &\text{Cor.~\ref{cor:C} + H3}
    \Rightarrow\text{Lem.~\ref{lem:E}}
    \Rightarrow\text{Cor.~\ref{cor:closure}}.
\end{align*}
\textbf{Key point:}
Lemma~\ref{lem:F} requires only Assumption~\ref{ass:gap}.
Proposition~\ref{prop:U} enters exclusively via
conditions~(b) and Term~2 of Corollary~\ref{cor:C},
and as one sufficient route to Assumption~\ref{ass:gap}.
It is not the logical foundation of the
Dyadic Separation Principle.
\end{remark}

\begin{remark}[Open problems]
\label{rem:open}
\begin{itemize}
\item \textbf{Assumption~\ref{ass:gap}:}
  a uniform gap lower bound; plausible from WKB
  or monotone comparison; no proof available.
\item \textbf{Proposition~\ref{prop:U}:}
  uniform WKB-to-Airy matching with
  $O(c^{-2/3})$ error; no proof available.
\item \textbf{B-strong:}
  $P_{kl}\le C_2 c^{1/2}$ in the transition zone;
  no proof available.
\item \textbf{Conjecture~\ref{conj:ULW}:}
  analytic derivation of the erfc law and
  determination of $\lim_{c\to\infty}\beta(c)$.
\end{itemize}
\end{remark}

% ============================================================
\begin{thebibliography}{9}
\bibitem{DLMF}
{\em NIST Digital Library of Mathematical Functions},
F.~W.~J.~Olver et al.\ (eds.),
Release 1.1.12, \url{https://dlmf.nist.gov/}, 2023.

\bibitem{paperVI}
[Author], \textit{Galerkin Norm Estimates for the Prolate
Phase Operator}, Paper VI in this series.

\bibitem{paperVII}
[Author], \textit{Dyadic Cancellation for PSWF Galerkin
Operators}, Paper VII in this series.

\bibitem{OsipovRokhlinXiao}
A.~Osipov, V.~Rokhlin, H.~Xiao,
\textit{Prolate Spheroidal Wave Functions of Order Zero},
Springer, 2013.
\end{thebibliography}

\end{document}
